\begin{figure}[tb]
\centering
\begin{tikzpicture}[
  font=\small,
  possible/.style={-{Stealth[length=1.7mm]}, dashed, semithick},
  flow/.style={-{Stealth[length=1.7mm]}, semithick},
  observed/.style={draw, fill=black!7, minimum height=0.95cm,
                   text width=1.85cm, align=center, inner sep=4pt},
  candidate/.style={draw, dashed, minimum height=0.9cm,
                    text width=3.5cm, align=center, inner sep=4pt}
]
\node[observed] (instruction) at (0,0) {Retained\\instruction};
\node[candidate] (policy) at (4.6,1.08)
  {Ordinary task state\\or answer policy};
\node[candidate] (self) at (4.6,-1.08)
  {Self-referential computation\\induced anew};
\node[align=center, font=\footnotesize\itshape] at (4.6,0)
  {These may coexist.};
\node[observed] (response) at (9.2,0) {Generated\\response};
\node[observed] (label) at (12.8,0) {Report\\label};
\draw[possible] (instruction.east)--(1.8,0)|-(policy.west);
\draw[possible] (1.8,0)|-(self.west);
\draw[possible] (policy.east)--(7.15,1.08)|-(response.west);
\draw[possible] (self.east)--(7.15,-1.08)|-(response.west);
\draw[flow] (response)--node[above, font=\footnotesize] {scoring}(label);
\end{tikzpicture}
\caption{\textbf{An instruction effect does not choose between these explanations.}
The experiment manipulates the retained instruction and measures response
labels. Dashed boxes and paths are compatible internal explanations, not
measured mediators or an exhaustive list. Ordinary task control and
self-referential computation may overlap or occur together. The solid final
arrow is the observed scoring operation, not validation of the report's
truth. Transcript and query are omitted here to isolate the unresolved
interpretation of the retained-instruction effect.}
\label{fig:compatible-explanations}
\end{figure}
